\documentclass[11pt]{article}
\usepackage[left=1in,top=0.8in,right=1in,bottom=0.8in]{geometry}
\usepackage[hidelinks]{hyperref}
\usepackage{mathptmx}
\usepackage{parskip}
\pagestyle{empty}

\begin{document}

{\Large\textbf{Rex Tin Long Fung}}\\
E-mail: rexfung@umich.edu \quad Website: \href{https://rextlfung.github.io/}{rextlfung.github.io}

\medskip
October 11, 2026

\medskip
Procter \& Gamble\\
Mason Business \& Innovation Center\\
Mason, OH

\medskip
\textbf{Re: R\&D PhD Summer Intern, Imaging Technologies (Requisition R000160079)}

Dear Hiring Team,

I am a PhD candidate in Biomedical Engineering at the University of Michigan, and I am writing to apply for the Imaging Technologies internship for summer 2027. My dissertation work is about building imaging systems that capture better data and then proving, with quantitative evidence, that they do. That is the work this role describes, and I would like to bring it to P&G's before-and-after imaging programs.

\textbf{Hardware and acquisition.} In the Functional MRI Lab I design and program the pulse sequences that control scanner hardware. I developed a randomly undersampled 3D-EPI acquisition that accelerates fMRI by $10\times$ and enables sub-2 mm resolution at 3T,. I also built open-source Pulseq sequences that run across scanner vendors and are now used by 7+ institutions, including Harvard and Freiburg. Making one capture workflow behave consistently on different hardware is the same problem as integrating and evaluating components that affect image capture.

\textbf{Image analysis and computer vision.} I reconstruct 100 GB+ datasets with custom iterative algorithms in Julia and Python, and I have written OpenCV pipelines for automated cell counting in 200+ microscopy images and behavioral tracking in 30+ hours of video, cutting manual analysis time by 90\%. Outside the lab, I fine-tuned a YOLOv8 detector for my Volleybot project that beat a pretrained baseline by 17 percentage points while running 6$\times$ faster. I am comfortable moving between Python and MATLAB, and between classical image processing and deep-learning workflows.

\textbf{Experimental rigor and repeatability.} I contributed to a multi-site protocol harmonization study evaluated with ABCD quality-control metrics, which made measurement repeatability and careful attention to experimental data a routine part of my work. I have less hands-on experience with optics, illumination, and motion-control hardware, and I would welcome the chance to learn them from P\&G's imaging experts. My background in embedded control systems and a habit of validating every change quantitatively should help me ramp up quickly.

\textbf{Communication.} I have a first-author manuscript under revision at \textit{Magnetic Resonance in Medicine}, co-authored a journal article in \textit{IEEE Transactions on Computational Imaging}, and presented at ISMRM, OHBM, and a Gordon Research Conference. As a teaching assistant for Medical Imaging Systems, I explained imaging physics to 20--30 students, which taught me to document and present technical work to people outside my specialty.

I am available for the full 12-week internship from May or early June through August 2027 and am prepared to work on site at the Mason Business \& Innovation Center four to five days a week. Thank you for your consideration; I would be glad to discuss how I can contribute.

Sincerely,

Rex Fung
\end{document}
